\section{Problema e cenário escolhido}

Entre os quatro cenários propostos pelo professor, o grupo escolheu o de
\textbf{processamento distribuído de grande volume de dados com arquitetura
mestre/workers}. O problema é processar um conjunto de dados grande demais
para ser tratado por uma única máquina em tempo razoável, dividindo o
trabalho entre várias máquinas que processam partes dos dados ao mesmo
tempo e depois combinam seus resultados.

\paragraph{Objetivo.} Construir um sistema mestre/workers que distribua o
processamento de um grande volume de dados entre várias máquinas e
avaliar experimentalmente seu desempenho (\emph{speedup}), o balanceamento
de carga entre máquinas heterogêneas e a tolerância a falhas de workers.

Este pré-projeto tem caráter \textbf{exploratório} em duas frentes: o
\textbf{conjunto de dados} (Seção~\ref{sec:datasets}), ainda em análise,
e o \textbf{ambiente de execução} (Seção~\ref{sec:ambiente}), cujas opções
foram avaliadas e que já está definido. A pergunta que o sistema vai
responder e o tipo de processamento (agregação, contagem, filtragem\ldots)
são decisões posteriores, tomadas após a definição do dataset.
